\newpage
\phantomsection
\label{prf:negation-of-product}

\begin{remark*}[Return]
\hyperref[thm:negation-of-product]{Return to Theorem}
\end{remark*}

\begin{theorem*}[Negation of a Product]
Let $F$ be a field and let $a,b\in F$.  Then
\[
-(a\cdot b) = (-a)\cdot b = a\cdot(-b).
\]
\end{theorem*}

\begin{proof}
\textbf{Professional Standard Proof.}~\\
We show $(-a)\cdot b$ is the additive inverse of $a\cdot b$ by computing
their sum:
\[
a\cdot b + (-a)\cdot b = (a + (-a))\cdot b = 0\cdot b = 0.
\]
By uniqueness of additive inverses, $(-a)\cdot b = -(a\cdot b)$.
The identity $a\cdot(-b) = -(a\cdot b)$ follows symmetrically:
\[
a\cdot b + a\cdot(-b) = a\cdot(b + (-b)) = a\cdot 0 = 0.
\]
\end{proof}

\begin{proof}
\textbf{Detailed Learning Proof.}~\\

We prove both equalities separately.

\textbf{Part A: $(-a)\cdot b = -(a\cdot b)$.}

\textbf{Step A1.} Compute $a\cdot b + (-a)\cdot b$.

By distributivity (right),
\[
a\cdot b + (-a)\cdot b = (a + (-a))\cdot b.
\]

\textbf{Step A2.} Apply the additive inverse axiom.

By the additive inverse axiom, $a + (-a) = 0$.  Therefore,
\[
(a + (-a))\cdot b = 0\cdot b.
\]

\textbf{Step A3.} Apply the zero product theorem.

By \hyperref[thm:zero-product]{Theorem~\ref*{thm:zero-product}} (with the
roles of $a$ and $b$ swapped, using commutativity of multiplication),
$0\cdot b = 0$.  Therefore,
\[
a\cdot b + (-a)\cdot b = 0.
\]

\textbf{Step A4.} Conclude by uniqueness of additive inverses.

The equation $a\cdot b + (-a)\cdot b = 0$ states that $(-a)\cdot b$ is an
additive inverse of $a\cdot b$.  By additive cancellation, additive inverses
are unique.  Therefore,
\[
(-a)\cdot b = -(a\cdot b).
\]

\textbf{Part B: $a\cdot(-b) = -(a\cdot b)$.}

\textbf{Step B1.} Compute $a\cdot b + a\cdot(-b)$.

By distributivity (left),
\[
a\cdot b + a\cdot(-b) = a\cdot(b + (-b)).
\]

\textbf{Step B2.} Apply the additive inverse axiom.

By the additive inverse axiom, $b + (-b) = 0$.  Therefore,
\[
a\cdot(b + (-b)) = a\cdot 0.
\]

\textbf{Step B3.} Apply the zero product theorem.

By \hyperref[thm:zero-product]{Theorem~\ref*{thm:zero-product}},
$a\cdot 0 = 0$.  Therefore,
\[
a\cdot b + a\cdot(-b) = 0.
\]

\textbf{Step B4.} Conclude by uniqueness of additive inverses.

The equation $a\cdot b + a\cdot(-b) = 0$ states that $a\cdot(-b)$ is an
additive inverse of $a\cdot b$.  By additive cancellation, additive inverses
are unique.  Therefore,
\[
a\cdot(-b) = -(a\cdot b).
\]

\textbf{Conclusion.}

Combining Parts A and B:
\[
-(a\cdot b) = (-a)\cdot b = a\cdot(-b). \qedhere
\]
\end{proof}

\begin{remark*}[Proof structure]
Both identities use the same strategy: show that the candidate equals
$-(a\cdot b)$ by verifying it sums with $a\cdot b$ to give $0$, then invoke
uniqueness of additive inverses via cancellation.  The computation in each
case is one application of distributivity, one application of the additive
inverse axiom, and one application of the zero product theorem.  Part A uses
right-distributivity; Part B uses left-distributivity.
\end{remark*}

\begin{remark*}[Dependencies]~\\
\begin{itemize}
  \item \hyperref[def:field]{Definition~\ref*{def:field}}
        \textnormal{(left and right distributivity; additive inverse axiom;
        commutativity of multiplication)}
  \item \hyperref[thm:zero-product]{Theorem~\ref*{thm:zero-product}}
        \textnormal{($a\cdot 0 = 0$ and $0\cdot b = 0$)}
  \item \hyperref[thm:additive-cancellation]{Theorem~\ref*{thm:additive-cancellation}}
        \textnormal{(uniqueness of additive inverses)}
\end{itemize}
\end{remark*}
